\begin{enumerate}[label=\textbf{\ansnum.\arabic*}]
  \item \verb!git config --global user.name "Ada Student"! and \verb!git config --global user.email "ada@example.org"!.
  \item Working directory, staging area, repository. \texttt{git add} moves a change from the working directory to the staging area, \texttt{git commit} from the staging area into the repository.
  \item \texttt{\ M walk.py} (modified, not staged) and \texttt{?? notes.py} (untracked).
  \item \texttt{git diff} compares the working directory with the staging area (changes not yet staged). \texttt{git diff --staged} compares the staging area with the last commit: exactly what the next commit will contain.
  \item \verb!git log -p config.py! (add \texttt{--oneline} or \texttt{--stat} for a shorter view).
  \item
\begin{verbatim}
__pycache__/
*.pyc
.venv/
*.log
!keep.log
.env
\end{verbatim}
  The exception must come \emph{after} the pattern it overrides.
  \item \verb!git rm --cached debug.log!, then commit. The file stays on disk and is ignored from now on.
  \item The grandparent of the newest commit (two steps back). Show it with \verb!git show HEAD~2!.
  \item \verb!git blame controller.py! shows the commit for every line; then \verb!git show COMMIT! (or \texttt{git log -p}) for the message and the context.
  \item[\textbf{\ansnum.10}] A commit is a permanent snapshot. The later commit only records the deletion; the earlier commit with the password is still part of the history and everyone who clones the repository can read it with \texttt{git log -p}. The password must be treated as published: \emph{change the password} first, and only then (optionally) rewrite the history with special tools.
\end{enumerate}
